\section{Roman}

The Romans took the Greek scheme almost whole. They matched their own gods to the
Greek gods and read the same theogony. The Latin retelling is \textbf{Ovid's
Metamorphoses}, which opens with \textbf{Chaos} resolved into the ordered world.
The distinct Roman contribution is in the \textbf{household and state religion}
rather than in the cosmogony.

\subsection{The Creation Model}

In the beginning there was a \textbf{formless mass}, a rough and unordered heap
in which the elements warred against one another. Earth, water, air, and fire were
jumbled together, each opposing the rest. There was no shape and no settled place
for anything.

Then a \textbf{shaping power} ended the conflict. It separated the earth from the
sky, the sea from the land, and the heavy from the light. It rounded the earth,
set the seas in their basins, lifted the air, and placed the stars. Last of all
came \textbf{man}, made to look up at the sky. The warring heap became an ordered
world with its realms laid out in order.

After this came the \textbf{ages of mortals}, a long decline. First a
\textbf{Golden Age} of ease and peace with no toil and no law. Then a
\textbf{Silver Age}, then a \textbf{Bronze Age}, and at last a hard \textbf{Iron
Age} of labor, greed, and war. The gold of the first age gave way step by step to
the iron of the present one.

\subsection{The Model with Its Names}

Where Hesiod begins with a pure genealogy, Ovid begins with a craftsman-like
\textbf{ordering power}. He does not name one supreme creator with certainty. He
writes only \textbf{whichever of the gods it was} who shaped the world. This Roman
account leans on the Greek throughout and adds little to the cosmogony itself.

\subsubsection{Janus and the Doorway}

\textbf{Janus} is a Roman god with no Greek match. He is the god of
\textbf{doorways}, \textbf{beginnings}, and \textbf{transitions}, two-faced and
looking both ways at once. In some accounts he is the \textbf{first god} of the
early world, present at the very start of things. His name opens the year, and
every Roman prayer began with him.

\subsubsection{The Greek-to-Roman Mapping}

The Roman gods are the Greek gods under Latin names, with a few shifts of
emphasis.

\begin{longtable}{L{0.22\linewidth} L{0.26\linewidth} L{0.30\linewidth}}
\caption{The Greek gods under their Roman names.}\\
\toprule
\textbf{Greek} & \textbf{Roman} & \textbf{Domain} \\
\midrule
Zeus & Jupiter & Sky, king of gods \\
Hera & Juno & Marriage, queen \\
Poseidon & Neptune & Sea \\
Hades & Pluto, Dis & Underworld \\
Demeter & Ceres & Grain, harvest \\
Aphrodite & Venus & Love, beauty \\
Ares & Mars & War \\
Athena & Minerva & Wisdom, craft \\
Hermes & Mercury & Messages, trade \\
Hephaestus & Vulcan & Fire, the forge \\
Artemis & Diana & Hunt, the moon \\
Apollo & Apollo & Sun, prophecy, healing \\
Dionysus & Bacchus, Liber & Wine, ecstasy \\
Kronos & Saturn & The deposed elder, the golden age \\
Persephone & Proserpina & Queen of the dead \\
\bottomrule
\end{longtable}

\subsubsection{The Household and State Religion}

Beside the great gods the Romans worshipped many small powers in and around
the home. A \textbf{numen} is a \textbf{divine presence} felt in a place, an act,
or a thing, not always a named god with a body. The early Roman sense of the holy
was this felt power more than the Greek personal god. A doorway, a boundary stone,
or a spring could each hold its numen.

\begin{longtable}{L{0.20\linewidth} L{0.54\linewidth}}
\caption{The household gods, worshipped at the home shrine, the lararium.}\\
\toprule
\textbf{Spirit} & \textbf{What it guards} \\
\midrule
Lares & The household, the family, the land and its boundaries \\
Penates & The storehouse, the food, the inner home \\
Genius & The spirit of the head of the house and the male line \\
Vesta & The hearth fire itself, kept ever burning \\
\bottomrule
\end{longtable}

The \textbf{Penates} of the Roman state were said to be those brought from
\textbf{Troy by Aeneas}, linking the household cult to the founding of Rome.

\subsubsection{The Underworld and the Cult of the Dead}

The Roman underworld is the Greek one under \textbf{Pluto} and \textbf{Proserpina},
with the same rivers, ferryman, and judges. The Romans honored their own dead
more closely as gods of a kind.

\begin{longtable}{L{0.22\linewidth} L{0.52\linewidth}}
\caption{The Roman cult of the dead.}\\
\toprule
\textbf{Spirit} & \textbf{Who they are} \\
\midrule
Manes & The spirits of dead ancestors, honored as kindly \\
Lemures & The restless or hostile dead, ghosts to be appeased \\
Lemuria & The festival in May to drive the Lemures from the house \\
Parentalia & The festival in February to honor the family dead \\
\bottomrule
\end{longtable}

The dead were fed and remembered at the tomb. The family line, living and dead,
was one continuing household under its gods.

\subsection{Comparison}

The Roman cosmology states the pattern of this collection in a borrowed form. The
\textbf{one goes out into the many and returns}. The going-out is the formless
mass resolved by a shaping power into the many ordered realms, then filled with
gods and mortals through the long ages. This is the same arc as the Greek, since
the Romans read the same theogony. Rome adds its own contribution at the near end
of the cycle, in the \textbf{household and the cult of the dead}, where each
family keeps its own small return. The living tend the spirits of their own dead,
and the line of ancestors carries on as one unbroken household. The base models
name a single source and map its full descent and repair. The Roman account,
leaning on the Greek, leaves the source as Ovid's uncertain \textbf{whichever of
the gods it was}. The same structure holds. The many are resolved out of a
formless first state, and the dead are gathered back into the continuing household
they came from.
